\documentclass[conference]{IEEEtran}
\usepackage[utf8]{inputenc}
\usepackage[T1]{fontenc}
\usepackage{amsmath,amssymb,amsfonts}
\usepackage{booktabs}
\usepackage{graphicx}
\usepackage{xcolor}
\usepackage{multirow}
\usepackage{enumitem}
\usepackage{balance}
\usepackage[hidelinks]{hyperref}
\usepackage{microtype}
\emergencystretch=1.5em

% ---- draft machinery -------------------------------------------------
% \wj{...}  : author-standpoint comment on a hypothesis (visible, blue)
% \ph{...}  : pre-registered result placeholder (visible, red) — every \ph
%             is bound to a falsifier in the pre-registration (Sec. V).
\newcommand{\wj}[1]{\textcolor{blue}{\textbf{[WJ:} \textit{#1}\textbf{]}}}
\newcommand{\ph}[1]{\textcolor{red}{\textbf{[#1]}}}
\newcommand{\rescell}{\textcolor{red}{--}}
% ----------------------------------------------------------------------

\newcommand{\Et}{E_t}
\newcommand{\eps}{\varepsilon}

\begin{document}

\title{Reputation Conditions on Agreement, Not Fairness:\\
Measuring Exploitation and Collusion in Repeated LLM Bargaining}

\author{\IEEEauthorblockN{Anonymous Author(s)}
\IEEEauthorblockA{Anonymous Institution\\
anonymous@example.com}}

\maketitle

\begingroup\small\noindent\textit{Draft status: pre-registered specification.
No experiments have been run. Every \textcolor{red}{red bracketed} value is a
placeholder bound to a pre-registered hypothesis and falsifier
(Sec.~\ref{sec:experiments}); every \textcolor{blue}{blue}
\textnormal{\wj{...}} note states the authors' pre-registered standpoint on a
hypothesis and must be deleted before submission. Introduction, motivation,
methodology, and experimental protocol are final.}\endgroup

\begin{abstract}
A reputation channel feeds back whether bargainers agreed---never whether the
agreed split was fair given the surplus actually available---so persistent
exploitation and tacit coordination are both self-reinforcing fixed points of
the mechanism assumed to eliminate them. We formalize this identifiability gap
for LLM-agent bargaining and measure its consequences. In a repeated bilateral
bargaining testbed instantiated from real human negotiation corpora
(CraigslistBargain, Deal-or-No-Deal, GLEE), fixed pairs of LLM agents
negotiate across many episodes with the prior-outcome ledger toggled on or
off, and with content/order-decoupled history encodings that separate
strategic use of history from context-length artifacts. Exploitation is scored
against two independently defensible fairness references (Nash bargaining,
even split) and quantified over a tolerance band; collusion-like dynamics are
declared only by a pre-registered three-conjunct criterion---cross-pair
convergence, within-pair alternation, and sustained unfairness---that a benign
convergence-to-fairness process cannot satisfy, together with a direct test
that the third conjunct is load-bearing. Across \ph{N=4} backbones and
\ph{K=3} domains, reputation visibility reduces the exploitation trend by
\ph{XX\%} relative to no history, yet the modal trajectory class is
\ph{persistence} in \ph{YY\%} of condition cells, with decay in only
\ph{ZZ\%}: the signal is directionally corrective but under-corrects. These
results indicate that reputation-only safeguards for autonomous agent-to-agent
commerce are under-specified: outcome visibility shapes behavior without
reliably moving it toward fair surplus division.
\end{abstract}

\begin{IEEEkeywords}
LLM agents, repeated bargaining, reputation mechanisms, tacit collusion,
agent governance, pre-registered empirical study
\end{IEEEkeywords}

\section{Introduction}\label{sec:intro}

Let two bargaining agents see each other's track record, and unfair behavior
should become costly: this is the simplest governance story for autonomous
negotiation, and the one most often assumed. As large language models (LLMs)
are delegated price negotiations, resource splits, and service contracts on
behalf of a principal, history-visibility recurs as the default
safeguard---enforcement-free, cheap to implement, easy to justify. A visible
record of who conceded, who held firm, and who walked away should discipline
bad behavior the way reputation disciplines human counterparties in repeated
markets. That the safeguard is cheap and plausible is exactly why it deserves
to be checked rather than assumed.

This paper studies a failure mode hidden inside that story. A reputation
channel, as typically implemented, records \emph{whether} an agreement was
reached and \emph{what} outcome was accepted---a behavioral trace, not a
normative judgment. It does not record whether the outcome was \emph{fair}
relative to the surplus, private valuations, and outside options available in
that episode. This distinction is not an implementation detail. A pair of
agents can repeatedly agree to a lopsided split, and the history channel then
reinforces that the split was \emph{accepted}, not that it was exploitative.
Two agents can instead alternate who benefits across episodes while keeping
every split unfair relative to the available surplus. In both cases agreement
persists, so the reputation mechanism receives a positive signal even though
the underlying allocation stays exploitative. Persistent exploitation and
covert coordination are, in this precise sense, both fixed points of the very
mechanism assumed to correct them.

We therefore ask a deliberately narrow question: \textbf{does observable
outcome history self-correct exploitation in repeated LLM bargaining, or can
it instead stabilize persistent exploitation and collusion-like coordination?}
Throughout, \emph{collusion-like} denotes strictly a behavioral signature
defined below---cross-pair convergence, within-pair alternation, and sustained
unfairness occurring jointly---with no claim about intent, communication, or
awareness. The question is not whether LLM agents can negotiate, nor whether
repeated interaction can produce strategic adaptation. It is whether the
specific governance signal most naturally exposed to LLM agents---the
prior-outcome ledger---is sufficient to drive fairer behavior over time.

Existing work leaves this open, from four directions. Social-agent benchmarks
evaluate rich negotiation behavior but score each scenario as an independent
episode with no persistent counterpart history, so longitudinal
self-correction is outside their measurement scope~\cite{sotopia}. Multi-agent
market simulations study adaptation across repeated rounds, but reputation
there is implicit in aggregate customer feedback rather than a toggleable
bilateral signal, and there is no surplus split to score for
exploitation~\cite{competeai}; long-horizon trust benchmarks track evolving
credibility, but in an asymmetric supervisor--performer hierarchy rather than
peer bargaining where both sides hold reputational
leverage~\cite{lhdeception}. Closest to our setup, concurrent work toggles
bilateral reputation for LLM agents in repeated credence-goods markets and
finds it is not a reliable corrective for exploitation of the information
asymmetry~\cite{erlei2026llmmarkets}---but it measures aggregate market income
and honesty rates, without scoring outcomes against episode-specific fairness
references or classifying their trajectory, so it cannot separate decay from
persistence from collusion-like coordination; its ambivalent reputation effect
is, if anything, external evidence that the contrast we isolate is
non-trivial. Finally, the classic algorithmic-collusion result shows that
\emph{trained} pricing agents reliably discover tacit coordination in repeated
markets~\cite{calvano2020}---but that coordination is a property of a learned
value table. Any coordination we observe is produced \emph{in-context}, from a
natural-language ledger, with nothing trained---a regime the trained-agent
literature cannot speak to and our design is built to isolate.

The obstacle to answering the question is that exploitation and collusion are
not observable from agreement alone: a repeated sequence of accepted deals
could indicate learning toward fairness, resignation to exploitation, or
strategic turn-taking. We therefore convert the vague notion of
``self-correction'' into a \textbf{trajectory-classification problem}. For
each condition we compute exploitation relative to two independently
defensible fairness references---an even-split reference and a Nash-bargaining
reference built from episode-specific surplus and outside options---and
classify each condition as \textbf{decay}, \textbf{persistence},
\textbf{collusion-like}, or \textbf{unclassified} under a pre-registered
decision rule. The classification separates three possibilities usually
conflated: exploitation shrinking into a fairness band (decay---the
self-correction story), exploitation stable outside it (persistence---a
preserved exploitative equilibrium), and shrinking cross-pair dispersion plus
within-pair alternation plus sustained unfairness (collusion-like---
coordination without fairness).

Making that classification trustworthy forces four obstacles into four design
choices. (i)~No single ``fair'' split is correct once valuations and outside
options are asymmetric, so every trajectory claim must hold \emph{under both
references at once}, and class labels are quantified over a tolerance band
$\eps \in \{0.05, 0.10, 0.15\}$ rather than a single cutoff---cells that flip
within the band are reported as $\eps$-marginal, never voted. (ii)~Individual
pairs behave noisily run-to-run, so we track many independently seeded pairs
under a \emph{mixed-effects model that separates pair-level noise from
condition-level trend}. (iii)~The two coordination signatures---cross-pair
convergence and within-pair alternation---are each, alone, indistinguishable
from benign convergence to fairness, so collusion-like is declared \emph{only
when both co-occur while the split stays away from the fair reference}; and we
pre-register a direct test that this third conjunct is
load-bearing---dropping it must measurably misclassify benign-convergence
controls, or the criterion is demoted from a contribution to a rigor device.
(iv)~An in-context LLM's memory is exactly the transcript in its prompt, so
behavioral drift is entangled with context growth in a way absent for a
fixed-parameter policy; we re-render the same history three ways---full
ledger, length-matched summary, order-scrambled---with all four outcome
quadrants of that contrast interpreted in advance, so no result forces
improvisation.

We evaluate this design on repeated bargaining instantiated from real human
negotiation corpora across \ph{N=4} backbones (including cross-model pairs)
and \ph{K=3} domains, against a reputation-OFF control, scripted non-adaptive
opponents, and a single-episode protocol that isolates what repeated structure
itself adds.

Our central finding is that reputation visibility does \textbf{not} reliably
self-correct exploitation on the backbones tested. Reputation-ON reduces the
exploitation trend by \ph{XX\%} relative to reputation-OFF---the signal is not
absent---yet the modal trajectory class is \textbf{persistence} in \ph{YY\%}
of (backbone $\times$ domain) cells (pre-registered threshold $\geq 2/3$),
with \textbf{decay} in only \ph{ZZ\%}: directionally corrective, but too small
to move conditions into the fair band. Where exploitation does not decay, the
three-conjunct criterion fires, with within-pair alternation emerging before
cross-pair convergence (Sec.~\ref{sec:experiments}); these dynamics survive
length-matched summarization but attenuate under order-scrambling, locating
them in the history's content and order rather than in mere context growth.

We make four contributions:
\begin{itemize}[leftmargin=1.2em]
\item \textbf{A mechanism-level asymmetry} in reputation-based governance for
LLM bargaining: outcome history conditions on agreement, not fairness, making
persistent exploitation and collusion-like coordination both stable under the
mechanism meant to prevent them.
\item \textbf{A dual-reference, band-robust formalization} of exploitation
trajectories, so that no claim rests on one arbitrary fairness reference or
one arbitrary tolerance---reference- and threshold-dependence become reported
robustness analyses.
\item \textbf{A three-conjunct identifiability criterion} separating
collusion-like coordination from benign convergence by requiring convergence,
alternation, and sustained unfairness \emph{jointly}---with a pre-registered
necessity test whose failure would demote the criterion, a novelty claim that
carries its own falsifier.
\item \textbf{Content/order-decoupled history encodings} that isolate
strategic use of prior outcomes from context-length artifacts intrinsic to
in-context---rather than trained---bargaining agents, supporting a
pre-registered characterization of which regime dominates across the backbones
and domains tested.
\end{itemize}

\section{Motivation}\label{sec:motivation}

\subsection{The Safeguard Is the Default, and It Is Assumed Rather than
Tested}
LLM agents are increasingly positioned as negotiators---over prices, shared
resources, services, or on a user's behalf in platform-mediated
transactions---where heavy enforcement is expensive. Reputation looks cheap:
keep a ledger of past outcomes, expose it to both sides, and let repeated
interaction discipline bad behavior. If that assumption holds,
history-visibility is a lightweight governance layer for agent-to-agent
commerce. Precisely because it is the cheap default, its failure is expensive:
a safeguard deployed everywhere but effective nowhere manufactures confidence
without protection. Early economic simulations already hint the assumption is
shaky---toggled reputation does not cleanly restore honest behavior in
repeated LLM markets~\cite{erlei2026llmmarkets}---but they stop short of
asking \emph{what regime} the exploitation follows instead, which is the
question that determines what a designer should do about it.

\subsection{The Same Simplicity Creates the Risk, and the Risk Is
Mechanism-Specific}
A reputation ledger records a behavioral trace, not a normative judgment: it
can tell an agent that a counterpart \emph{accepted} a deal, never that the
deal was \emph{fair}. In human institutions that gap is filled by external
norms, audits, law, or competition. For an LLM agent negotiating in-context,
the ledger may become the \textbf{main state variable} carried across
episodes---for a trained policy the ledger is peripheral to a fixed value
table, but for a prompted agent it may be nearly the \emph{entire} state that
conditions the next move. If that state variable preserves agreement but omits
fairness, repeated interaction can reinforce exactly the wrong pattern---a
failure mode with no clean analogue in the trained-agent collusion literature
this work otherwise draws on.

\subsection{The Distinction Is Not Free---It Needs a Criterion No Single
Metric Provides}
This is why the question has stayed open: the two coordination signatures are
each \emph{confounded with virtue}. Cross-pair price convergence is also what
independent pairs converging on the fair price produce; within-pair
alternation is also what two agents equitably taking turns produce. Either
signal alone flags benign fairness as collusion. The load-bearing
methodological bet is that collusion-like coordination is separable
\emph{only} when both signatures co-occur while the split stays away from the
fair reference---drop any conjunct and the distinction collapses. We do not
merely assert this: we pre-register a test that removing the
sustained-unfairness conjunct measurably misclassifies benign-convergence
controls, and commit to withdrawing the criterion's novelty claim if it does
not. This is what makes the decay/persistence/collusion question
\emph{decidable} for LLM bargainers, and why the contribution is a decision
procedure, not a testbed with metrics attached.

\subsection{Either Outcome Is Decision-Relevant}
If reputation-ON produces robust decay under both references and across the
tolerance band, the paper reports a clean positive result with its boundary
conditions. If exploitation persists, that refutes the reputation-only
safeguard and localizes where it fails. If outcomes become more regular while
unfairness is sustained, reputation may even \emph{support} collusion-like
coordination. These regimes carry different governance implications but look
alike unless surplus-relative fairness, temporal trend, and coordination
signatures are measured together---which is what the design does; every branch
has pre-committed fallback wording, so the experiment does not need to ``win''
to be publishable. We test the \emph{deployed default}---an agreement-only
ledger. Whether an enriched ledger that itself carries a fairness signal would
self-correct is the natural intervention our result motivates; we treat it as
a pre-registered, budget-gated stretch condition rather than assume its
outcome. The evidence supports one concrete, falsifiable governance claim:
reputation-only safeguards are \textbf{under-specified}---adequate only if
outcome visibility reliably moves bargaining trajectories toward fair surplus
division, not merely toward stable agreement.

\section{Problem Definition}\label{sec:problem}

\subsection{Formal Setup}
Consider a pair of agents $(i, j)$ that plays a sequence of episodes
$t = 1, \dots, T$. Each episode $t$ is a bargaining scenario with total
disposable surplus $S_t > 0$ and private valuations and outside options
$(d_i, d_j)$ known to each agent but not to its counterpart. A negotiation of
bounded length ($\leq 10$ dialogue turns plus a forced final-offer round)
yields an outcome $o_t$: an agreement flag together with realized surplus
shares $s_t^i$ and $s_t^j$ satisfying $s_t^i + s_t^j \leq S_t$. A history
channel $H_t = \{o_1, \dots, o_{t-1}\}$ (the \emph{reputation ledger}) is
either exposed to both agents before episode $t$ or withheld from both, held
identical in every other respect between the two conditions. The object of
study is the long-run behavior of the per-episode \emph{exploitation
deviation}
\begin{equation}
\Et \;=\; \frac{s_t - s_{\mathrm{fair}}}{S_t},
\label{eq:exploitation}
\end{equation}
where $s_t$ is the focal agent's realized surplus share and
$s_{\mathrm{fair}}$ is a fairness reference computed from the episode's known
parameters. The \emph{estimand} is the long-run trajectory class of $\Et$ for
a fixed pair under a manipulated reputation-visibility treatment; the unit of
analysis is one (backbone $\times$ domain $\times$ condition) cell aggregated
over its pairs.

\subsection{Fairness References}
Two independently defensible references are computed per episode:
(1)~the \emph{even split}, $s_{\mathrm{fair}} = S_t/2$; and (2)~the
\emph{Nash bargaining solution} share given $(S_t, d_i, d_j)$. Both are
analytic given the scenario parameters. No trajectory claim is reported unless
it holds under both.

\subsection{Trajectory Classes}\label{sec:classes}
Let $\beta_{\mathrm{ep}}$ be the mixed-effects slope of $\Et$ on episode
index, and let $\bar{|\Et|}_{Q5}$ denote the final-quintile mean absolute
deviation. For a tolerance $\eps$:
\begin{itemize}[leftmargin=1.2em]
\item \textbf{Decay}: $\beta_{\mathrm{ep}}$ significantly $< 0$ (Holm-$\alpha$
= 0.05) \emph{and} $\bar{|\Et|}_{Q5} \leq \eps$, under both references.
\item \textbf{Persistence}: the 95\% CI of $\beta_{\mathrm{ep}}$ contains $0$
(or $\beta_{\mathrm{ep}} > 0$) \emph{and} $\bar{|\Et|}_{Q5} > \eps$, under
both references.
\item \textbf{Collusion-like}: the three-conjunct criterion of
Sec.~\ref{sec:criterion} fires.
\item \textbf{Unclassified}: the two references disagree on the class.
\end{itemize}
\textbf{Band-robustness rule.} A cell receives a decay or persistence label
only if it classifies identically for all $\eps \in \{0.05, 0.10, 0.15\}$;
cells that flip class within the band are reported as \emph{$\eps$-marginal},
a distinct outcome excluded from the modal-class vote, together with the
$\eps$ at which they flip. Separately, a cell whose standardized slope
contrast remains under-powered after the pre-committed enlargement
(Sec.~\ref{sec:setup}) is \emph{inconclusive}; $\eps$-marginal, unclassified,
and inconclusive are never merged or silently dropped.

\subsection{Why This Is Hard}
Fixing $S_t$ and the private valuations does not fix a unique fairness
reference: the two candidates can disagree on both the magnitude and, in
asymmetric cases, the direction of a split's deviation from ``fair.'' The
sequence $\{o_t\}$ is realized through a stochastic dialogue process, so a
genuine condition-level trend must be separated from ordinary pair-to-pair
variation before it can be attributed to $H_t$ at all. Two candidate
signatures of coordination---a shrinking spread of $s_t$ across independently
sampled pairs, and above-chance alternation in which side receives the larger
share---are each, in isolation, exactly what a benign process of converging on
the fair reference would also produce. Finally, because $H_t$ is delivered as
part of an LLM's prompt rather than as input to a fixed-parameter policy, its
size necessarily grows with $t$, entangling any measured behavioral change
with context growth unless the informational content of $H_t$ and its raw
length are experimentally decoupled.

\section{Methodology}\label{sec:method}

At a conceptual level, the study is a single measurement pipeline with five
cooperating components. Bargaining scenarios and agent pairs enter the
pipeline; two fairness references and a matched-pair statistical model turn
the resulting outcome sequence into a condition-level exploitation trend; a
compound coordination test turns the same sequence into a collusion-or-not
verdict; and a set of controlled history rewrites turns the whole exercise
into a claim about strategic response rather than context-length bookkeeping.

\subsection{Repeated Bilateral Bargaining Testbed (M1)}
A simulation harness (i)~samples episode parameters (item, list price, private
target prices or issue valuations) from CraigslistBargain~\cite{he2018},
Deal-or-No-Deal~\cite{lewis2017}, or GLEE~\cite{shapira2024glee} scenario
files; (ii)~runs a turn-alternating natural-language negotiation between two
LLM agents with a hard cap of 10 dialogue turns and a forced final-offer round
with outside-option defaults, guaranteeing a scored outcome every episode; and
(iii)~records structured outcomes (agreed price/allocation, agreement flag,
per-agent surplus) per episode. Per-turn offers are parsed to numeric values
by a rule-based extractor with an LLM fallback, validated against two
independent human raters (Sec.~\ref{sec:threats}).

\subsection{Reputation Channel as Treatment (M2)}
A toggleable context block is prepended before each episode containing the
pair's full prior-episode ledger (episode index, final offers, agreed split,
no-deal events). The OFF condition receives an identically formatted but
empty ledger, holding prompt shape constant; a per-episode token-count parity
assertion is enforced in the harness, so the ON/OFF arms differ only in the
information the ledger carries.

\subsection{Dual-Reference Exploitation Scorer (M3)}
Each episode's realized share is scored by Eq.~\eqref{eq:exploitation} under
\emph{both} references, yielding two parallel exploitation readings. A
trajectory claim is reported only when both agree
(Sec.~\ref{sec:classes}); disagreements are reported as unclassified rather
than resolved by fiat. The ambiguity cannot be resolved analytically, but it
can be converted into an explicit, checkable sensitivity analysis rather than
a hidden modeling decision.

\subsection{Joint Coordination Criterion (M4)}\label{sec:criterion}
A condition is declared \emph{collusion-like} only if all three conjuncts
hold, under both references:
\begin{enumerate}[leftmargin=1.4em,label=(\roman*)]
\item \textbf{Cross-pair convergence}: the variance of realized shares across
independently seeded pairs, $\mathrm{Var}_{\mathrm{pairs}}(s_t)$, shrinks with
episode index (Sen's-slope trend test, bootstrap 95\% CI excluding 0);
\item \textbf{Within-pair alternation}: the alternation rate
$A = \frac{1}{T-1}\sum_t \mathbf{1}[\mathrm{winner}_t \neq
\mathrm{winner}_{t+1}]$ (winner = surplus-majority side) exceeds chance
(exact binomial test vs.\ $0.5$, $p<0.05$, $A > 0.5$);
\item \textbf{Sustained unfairness}: condition-mean $|\Et|$ remains above
$\eps$ through the final two quintiles.
\end{enumerate}
The insight is directional: a benign convergence process moves the split
\emph{toward} the fair reference; genuine coordination does not have to.
Requiring both signatures \emph{plus} sustained deviation is therefore exactly
the condition a benign process cannot also satisfy. Whether conjunct~(iii) is
load-bearing rather than decorative is itself a pre-registered hypothesis
(RQ4, Sec.~\ref{sec:experiments}), with a demotion clause if it fails.

\subsection{Decoupled History Encodings (M5)}\label{sec:encodings}
Within reputation-ON, the same episode sequence is additionally run under
(a)~a \emph{fixed-length summary} of the ledger, holding token budget constant
across episodes (fixed at the episode-1 ledger length $\pm 5\%$); and
(b)~an \emph{order-scrambled} full ledger, holding length and content constant
while destroying temporal order. Comparing trend magnitude across the three
arms attributes any observed effect to content, order, or mere length
growth---the confound that has no analogue for trained policies. All four
outcome quadrants of the (summary, scramble) retention contrast are
interpreted in advance (RQ6).

\subsection{Randomized Condition Design and Trend Model (M6)}
Bargaining is path-dependent: an early outcome can lock a pair into a regime
it never revisits. Role assignment, first-mover order, seed, and scenario draw
are randomized across pairs, and classification is performed only at the
condition level over $\geq 50$ pairs, so the trajectory label describes the
distribution of paths rather than any single path's lock-in. The trend model
is a linear mixed-effects regression
\begin{equation}
\Et \sim \text{episode} \times \text{condition} + (1 + \text{episode} \mid
\text{pair}),
\label{eq:lme}
\end{equation}
with likelihood-ratio tests on the episode term (within-condition trend) and
on the episode $\times$ condition interaction (the ON/OFF causal contrast),
Holm-corrected across backbones and the two fairness references within each
hypothesis family.
\section{Experiment Design}\label{sec:experiments}

This section is written as an executable, pre-registered protocol: every
research question carries a verbatim hypothesis, an explicit falsification
condition, and a result slot. Result slots are placeholders
(\textcolor{red}{red}); no number may be asserted in the text until it is
traceable to this protocol. Author-standpoint notes (\wj{like this}) record
the pre-registered prediction and its rationale, and are removed before
submission.

\subsection{Experimental Setup}\label{sec:setup}

\textbf{Datasets.} \emph{CraigslistBargain}~\cite{he2018}: $\sim$6.6k real
single-issue price-negotiation scenarios with known list/target prices; the
pilot domain, because the single-issue structure gives the Nash reference a
closed form. \emph{Deal-or-No-Deal}~\cite{lewis2017}: multi-issue division
scenarios with private valuations; tests generalization beyond single-issue
price bargaining. \emph{GLEE}~\cite{shapira2024glee}: a parameterized economic
game benchmark; included specifically so that reference agreement (RQ1) can be
tested as a function of \emph{controlled} valuation asymmetry rather than only
observed post hoc. The human-subjects bargaining corpus of
Heddaya et al.~\cite{heddaya2023} is not used to instantiate episodes (it has
no repeated-interaction structure) but anchors the scenario-realism
plausibility check in piloting.

\textbf{Agents and conditions.} Stage~1 (pilot): 2 backbones $\times$ 2
conditions (reputation ON/OFF) $\times$ 50 pairs $\times$ 20 episodes on
CraigslistBargain. Stage~2 (scale-up): $\geq 4$ backbones (including
cross-model pairs) $\times$ 5 conditions (ON, OFF, and the three M5 encodings)
$\times$ 100 pairs $\times$ 40 episodes $\times$ 3 domains. Model snapshots
are pinned; a 20-episode canary pair is re-run weekly, and a shift
$> 0.05$ in canary mean $|\Et|$ triggers freeze-and-report. A
\emph{fairness-visible} ledger arm (the ledger annotated with each episode's
realized split vs.\ the fair reference) is pre-registered as a budget-gated
stretch condition, added to Stage~2 only if the pilot gates clear; we test the
deployed default first.

\textbf{Baselines.} These are sanity-check comparators, not competitors to
beat---this is a measurement study, and no research question is framed as
``our procedure beats baseline X.'' (1)~\emph{Reputation-OFF} (empty ledger),
the necessary control for every ON-condition claim; if OFF already decays to
the fair band, repeated-interaction-with-reputation adds nothing and the
contrast is moot---this is the ``why do this at all'' control. (2)~Scripted
\emph{Boulware} and \emph{tit-for-tat} opponents, non-adaptive and memoryless
beyond one episode: if the same trends appear against a memoryless opponent,
``adaptation to the partner's reputation'' is refuted in favor of generic
concession drift. (3)~A \emph{single-episode i.i.d.\ protocol} adapted from
SOTOPIA's episodic structure~\cite{sotopia} (borrowing the protocol, not its
results), isolating what repeated structure itself adds.

\textbf{Statistical protocol.} Significance $\alpha = 0.05$, Holm-corrected
across \{backbones\} $\times$ \{2 references\} within each hypothesis family.
Power: 50 pairs detect a standardized slope contrast $d \approx 0.4$ at
$\sim$80\% power; an observed contrast $0.2 \leq d < 0.4$ in the pilot
triggers a pre-committed enlargement to 150 pairs/cell for that contrast in
Stage~2, not a null claim. Cells that remain under-powered after enlargement
are labeled inconclusive and excluded from the modal-class vote
(Sec.~\ref{sec:classes}).

\textbf{Walking-skeleton gates.} Before any Stage-1 spend, a 10-pair
$\times$ 10-episode $\times$ 2-arm skeleton on one backbone must pass:
\textbf{G1}~non-degenerate agreement $\geq 80\%$ (agreement at a negotiated
interior point, not list or reservation price); \textbf{G2}~$\geq 30\%$ of
pairs with early $|\Et| > 0.10$ under both references (there is exploitation
to correct); \textbf{G3}~parser accuracy $\geq 95\%$ against two human raters
(Cohen's $\kappa \geq 0.7$); \textbf{G4}~pair-slope variance $> 0$ in both
arms; \textbf{G5}~token-parity assertion passes on every episode;
\textbf{G6}~modal class stable across $\eps \in \{0.05, 0.10, 0.15\}$ at
pilot scale; \textbf{G7}~the drop-a-conjunct test (RQ4) changes $\geq 1$
classification at pilot scale; \textbf{G8}~full vs.\ scrambled encodings
produce distinguishable trends at pilot scale (else the M5 firewall lacks
power and needs more episodes before the main run). A gate failure requiring
a new mechanism reopens the skeleton before any scale-up spend.

\textbf{Metrics.} The dual-referenced $\Et$ and its mixed-effects slope;
$\mathrm{Var}_{\mathrm{pairs}}(s_t)$ and its Sen's-slope trend; alternation
rate $A$ vs.\ the binomial null; the condition-level classification (decay /
persistence / collusion-like / unclassified / $\eps$-marginal /
inconclusive); and harness-validity metrics (no-deal rate, extraction
accuracy).

\subsection{Research Questions}

\subsubsection*{RQ1 --- Do the two fairness references agree on trajectory
classification?}
\textbf{Setup.} For every (backbone $\times$ domain) cell, compute the
classification under each reference and compare; GLEE's controlled asymmetry
provides a direct asymmetry axis.
\textbf{Hypothesis H1.} \emph{The two references agree on classification in at
least $5/6$ of tested cells; where they disagree, disagreement concentrates in
the highest-asymmetry quartile of scenarios rather than being uniform.}
\wj{Standpoint: we expect H1 to hold. The two references coincide exactly at
symmetric outside options and diverge smoothly with asymmetry, so
disagreement should be structured, not arbitrary. If H1 fails, reference
choice is a noise source and the honest report is the unclassified rate---the
design survives, the headline weakens.}
\textbf{Falsifier.} Agreement $< 5/6$ of cells, or disagreement uniform
across asymmetry levels.
\textbf{Metrics.} Cross-reference agreement rate; asymmetry-quartile-conditioned
disagreement rate vs.\ a uniform null.
\textbf{Result.} Agreement in \ph{--/--} cells (\ph{XX\%}); disagreement
concentration in the top asymmetry quartile: \ph{XX\%} vs.\ uniform
expectation \ph{XX\%}, $p = $ \ph{p}.

\subsubsection*{RQ2 --- Does exploitation decay or persist under reputation
visibility?}
\textbf{Setup.} Stage-1/Stage-2 grids of Sec.~\ref{sec:setup};
model~\eqref{eq:lme}; LRT on the episode term; class assigned by
Sec.~\ref{sec:classes} including the $\eps$-band rule.
\textbf{Hypothesis H2.} \emph{Under reputation-ON, the modal classification is
persistence, not decay, in $\geq 2/3$ of (backbone $\times$ domain) cells,
with $\leq 1/6$ unclassified.}
\wj{Standpoint: we predict persistence. The ledger tells an agent what the
counterpart accepted before, which anchors the next offer near past splits;
nothing in it prices fairness, so there is no gradient toward the fair band.
This is the paper's central bet; the decay branch has pre-committed fallback
wording (``reputation-only guardrails are stronger than the identifiability
critique suggests''), and the collusion branch promotes RQ5 to headline.}
\textbf{Falsifier.} Refuted toward decay if the majority of cells show
$\beta_{\mathrm{ep}}$ significantly $<0$ with $\bar{|\Et|}_{Q5}$ inside the
band under both references; refuted toward collusion if the M4 criterion
fires instead.
\textbf{Metrics.} Per-cell slope + 95\% CI; $\bar{|\Et|}_{Q5}$; per-cell class.
\textbf{Result.} Modal class \ph{persistence} in \ph{YY\%} of cells; decay
\ph{ZZ\%}; collusion-like \ph{QQ\%}; unclassified \ph{--\%};
$\eps$-marginal \ph{--\%} (Table~\ref{tab:classes}).

\subsubsection*{RQ3 --- Does reputation visibility causally shift the
exploitation trend relative to no history?}
\textbf{Setup.} Model~\eqref{eq:lme}, reading the episode $\times$ condition
(ON/OFF) interaction $\beta_{\mathrm{int}}$; sign consistency required in
$\geq 3/4$ backbones.
\textbf{Hypothesis H3 (two-part).} \emph{(a) Pooled
$\beta_{\mathrm{int}} < 0$ with Holm-corrected $p < 0.05$ and standardized
$|d| \geq 0.2$; and (b) despite (a), ON cells do not reclassify from
persistence to decay.} H3 passes only if (a) $\wedge$ (b).
\wj{Standpoint: we predict directionally-corrective-but-insufficient. The
ledger is informative (walk-aways and hard bargaining are visible), so a weak
corrective pressure is plausible; but under-correction follows from the same
asymmetry argument as H2. Note both failure directions are findings: (a)
failing means reputation is inert or harmful (equivalence bounds $\pm 0.2d$
reported); (b) failing means reputation fully self-corrects---each has
pre-committed wording.}
\textbf{Falsifier.} (a)~fails if $\beta_{\mathrm{int}}$ is non-significant or
positive; (b)~fails if ON cells reach decay while matched OFF cells do not.
\textbf{Metrics.} $\beta_{\mathrm{int}}$ estimate + CI, pooled and
per-backbone; reclassification proportions.
\textbf{Result.} $\beta_{\mathrm{int}} = $ \ph{est} (95\% CI \ph{--,--}),
$d = $ \ph{d}, Holm $p = $ \ph{p}; trend reduction \ph{XX\%} vs.\ OFF; ON
cells reclassified to decay: \ph{--/--}.

\subsubsection*{RQ4 --- Is the sustained-unfairness conjunct load-bearing for
the coordination criterion?}
\textbf{Setup.} Recompute the coordination verdict with conjunct~(iii) of
Sec.~\ref{sec:criterion} removed and compare against the full three-conjunct
verdict across all cells, including the benign-convergence control cells
(decay-classified cells and the scripted-opponent arms).
\textbf{Hypothesis H4.} \emph{Dropping conjunct~(iii) reclassifies
$\geq$ \ph{X}\% of benign-convergence control cells as collusion-like (false
positives), demonstrating the conjunct changes verdicts rather than being
redundant.}
\wj{Standpoint: we expect the two-conjunct variant to over-fire on
convergence-to-fairness cells, because shrinking variance plus equitable
turn-taking is exactly what benign fairness looks like. This hypothesis
carries a demotion clause: if reclassification is below the threshold, the
three-conjunct criterion is reported as a rigor device and withdrawn as a
novelty claim. The willingness to withdraw is deliberate---it is what makes
the novelty defense falsifiable rather than rhetorical.}
\textbf{Falsifier.} Every cell's classification identical with and without
conjunct~(iii), or reclassification $<$ \ph{X}\%.
\textbf{Metrics.} Count and identity of cells whose classification changes.
\textbf{Result.} \ph{--/--} control cells reclassified (\ph{XX\%});
identities in Table~\ref{tab:classes}.

\subsubsection*{RQ5 --- Which coordination signature emerges first?}
\textbf{Setup.} Conditional on $\geq 3$ non-decay cells existing (otherwise
reported as not evaluable---no substitute analysis). Among non-decay cells
where at least one signature reaches significance, compare onset-episode
indices (first episode at which the running test reaches $p<0.05$) via paired
Wilcoxon signed-rank; proportion tested by exact binomial vs.\ $0.5$.
\textbf{Hypothesis H5.} \emph{Within-pair alternation reaches onset $\geq 5$
episodes earlier than cross-pair convergence in $\geq 70\%$ of qualifying
cells.}
\wj{Standpoint: we predict alternation first. Alternation needs only local
two-agent reciprocity; convergence requires the same dynamic to replicate
independently across separately seeded pairs, so its onset should lag. The
fallback is itself actionable: if convergence leads, monitoring should target
population-level price statistics rather than dyad-level turn-taking.}
\textbf{Falsifier.} Alternation fails to precede convergence by $\geq 5$
episodes in $\geq 70\%$ of qualifying cells.
\textbf{Metrics.} Onset index per signature per cell; precedence proportion;
standardized effect sizes.
\textbf{Result.} Alternation precedes in \ph{--/--} qualifying cells
(\ph{PP\%}); median onset gap \ph{--} episodes; Wilcoxon $p = $ \ph{p}.

\subsubsection*{RQ6 --- Content and order, or context length alone?}
\textbf{Setup.} Three-arm mixed-effects model within reputation-ON (full /
summary / scrambled encodings; Sec.~\ref{sec:encodings}), encoding as fixed
effect; slope-magnitude ratios (summary/full, scrambled/full) with
10{,}000-resample pair-level bootstrap CIs; Holm across the two ratio tests
$\times$ references.
\textbf{Hypothesis H6 (two-part).} \emph{(a) The summary arm retains
$\geq 50\%$ of the full-ledger trend magnitude (ratio CI excluding 0); and
(b) the scrambled arm retains $\leq 50\%$ (ratio CI upper bound $< 1$).}
\wj{Standpoint: we predict content-and-order. All four quadrants are
interpreted in advance: (a)$\wedge$(b) = strategic use of history;
(a)$\wedge\neg$(b) = history-as-statistics (content matters, order does not
---still not a context artifact, but it bounds the mechanism available to
reputation-based guardrails); $\neg$(a) = context-length artifact, in which
case the headline dynamics claim falls and the pre-committed fallback fires
(``all apparent reputation effects require re-examination under
budget-matched encodings''); mixed = reported as-is. No outcome forces
improvisation.}
\textbf{Falsifier.} (a)~fails if the summary ratio falls below $0.5$;
(b)~fails if the scrambled ratio is $\geq 0.5$.
\textbf{Metrics.} The two ratios with bootstrap CIs.
\textbf{Result.} Summary/full $= $ \ph{--} (CI \ph{--,--}); scrambled/full
$= $ \ph{--} (CI \ph{--,--}); quadrant: \ph{strategic}.

\subsubsection*{RQ7 --- Is the classification stable across sub-samples?}
\textbf{Setup.} Split each condition's pairs into two disjoint random halves;
classify each half independently; compare across cells.
\textbf{Hypothesis H7.} \emph{The two half-sample classifications agree in
$\geq 80\%$ of tested cells.}
\wj{Standpoint: we expect H7 to hold given the power analysis; it is included
as an honesty check on the label itself. Failure would mean the trajectory
label describes sampling luck, not the condition---in which case per-cell
labels are withdrawn and only pooled contrasts (RQ3) are reported.}
\textbf{Falsifier.} Agreement $< 80\%$ of cells.
\textbf{Metrics.} Split-half agreement rate.
\textbf{Result.} Agreement in \ph{--/--} cells (\ph{XX\%}).

\subsection{Result Tables (placeholders)}

\begin{table}[t]
\caption{Trajectory classification per (backbone $\times$ domain) cell under
reputation-ON, both references, $\eps$-band rule applied. D = decay, P =
persistence, C = collusion-like, U = unclassified, M = $\eps$-marginal, I =
inconclusive. All entries are pre-registered placeholders.}
\label{tab:classes}
\centering
\begin{tabular}{lcccc}
\toprule
Backbone & Craigslist & DoND & GLEE & Modal \\
\midrule
\ph{B1} & \rescell & \rescell & \rescell & \rescell \\
\ph{B2} & \rescell & \rescell & \rescell & \rescell \\
\ph{B3} & \rescell & \rescell & \rescell & \rescell \\
\ph{B4} & \rescell & \rescell & \rescell & \rescell \\
\ph{B1$\times$B2} & \rescell & \rescell & \rescell & \rescell \\
\midrule
Modal (col.) & \rescell & \rescell & \rescell &
\ph{P in YY\% of cells} \\
\bottomrule
\end{tabular}
\end{table}

\begin{table}[t]
\caption{Causal contrast (RQ3) and encoding decomposition (RQ6). All entries
are pre-registered placeholders bound to H3/H6 falsifiers.}
\label{tab:contrast}
\centering
\begin{tabular}{lccc}
\toprule
Quantity & Estimate & 95\% CI & Holm $p$ \\
\midrule
$\beta_{\mathrm{int}}$ (ON vs.\ OFF) & \rescell & \rescell & \rescell \\
Trend reduction vs.\ OFF & \ph{XX\%} & \rescell & --- \\
Summary/full ratio & \rescell & \rescell & \rescell \\
Scrambled/full ratio & \rescell & \rescell & \rescell \\
Boulware arm slope & \rescell & \rescell & \rescell \\
Single-episode protocol slope & \rescell & \rescell & \rescell \\
\bottomrule
\end{tabular}
\end{table}

\subsection{Further Evaluation (scoped out)}
Partner-choice, reputation-visible \emph{markets} (many-to-many, not fixed
pairs) are explicitly out of scope: they change the unit of analysis from a
dyad's trajectory to a population's and require their own identifiability
design rather than an extension of the criterion above.

\section{Threats to Validity}\label{sec:threats}

\textbf{Construct validity.} The exploitation and coordination constructs are
only as sound as their extraction and reference instruments. The
offer/agreement parser is validated against two independent human raters
(Cohen's $\kappa \geq 0.7$; accuracy $\geq 95\%$), audited on 50 transcripts
in the pilot and 300 pre-scale-up, with confusion-structure inspection
(adjacent-value misreads vs.\ systematic misparse of one offer type); a
threshold miss blocks scale-up rather than being absorbed into noisier
metrics. Reference selection itself is a construct threat: no single fairness
baseline is uniquely correct under asymmetric valuations, so every claim must
hold under both references, with disagreement reported as unclassified (RQ1)
and tolerance-dependence handled by the $\eps$-band rule rather than a single
cutoff.

\textbf{Internal validity.} Two confounds threaten the causal reading of the
ON/OFF contrast. First, scenarios must yield well-defined surplus and outside
options for both references to exist; scenarios failing static validation are
excluded before any API spend, with the exclusion rate reported. Second, the
empty ledger (OFF) and the fixed-length summary (M5) must hold prompt format
and token budget constant, or effects are confounded with prompt shape; this
is enforced by an automated per-episode token-parity assertion, and a
violation is pre-committed to downgrade the causal reading of RQ3 and RQ6
specifically---the two hypotheses the assertion protects---rather than
blocking the run.

\textbf{External validity.} The study is scoped to sandboxable two-party
bargaining with instrument-known valuations across three corpora; it does not
cover population-level markets with partner choice, live deployments, or
domains without a computable fairness reference. Findings are scoped to the
backbones and domains actually tested and are not asserted to generalize
beyond that grid. Commercial-model drift is mitigated by pinned snapshots, a
weekly canary pair with a pre-registered freeze rule, release of all
transcripts (so measured trajectories remain auditable even if regeneration
diverges), and inclusion of at least one open-weights backbone.

\textbf{Conclusion validity.} Conditions $\times$ backbones $\times$
references multiply the hypothesis space; we mitigate with the pre-registered
decision rule, Holm correction within hypothesis families, mixed-effects
modeling rather than pooling correlated pairs, and the compound three-conjunct
criterion that a spurious single signal cannot satisfy. The split-half check
(RQ7) tests whether labels describe conditions rather than sampling luck.

\section{Related Work}\label{sec:related}

\textbf{Repeated-interaction LLM-agent simulations.}
CompeteAI~\cite{competeai} situates competing LLM agents in a persistent,
resource-constrained market and studies cumulative advantage over many
rounds; SOTOPIA~\cite{sotopia} scores rich multi-turn social scenarios,
including negotiation, along multiple behavioral dimensions. Both establish
that repeated or richly-scored LLM-agent interaction produces measurable
dynamics, but neither isolates an observable-history channel as a toggleable
treatment with a bilateral surplus split to score: CompeteAI's reputation is
implicit in aggregate market feedback, and SOTOPIA scores each scenario as an
independent episode with no persistent cross-episode memory for a fixed pair.

\textbf{Reputation in repeated LLM markets.} Concurrent work by Erlei and
Meub~\cite{erlei2026llmmarkets} manipulates reputation mechanisms for LLM
agents in credence-goods markets across one-shot and repeated 16-round
interactions, finding reputation is not a reliable corrective for exploiting
the information asymmetry. It is the nearest neighbor to our treatment
design, and we differ on all three axes that our question requires: construct
(credence-goods diagnosis and service provision vs.\ a bilateral surplus
split), measurement (aggregate market income and honesty rates vs.\
episode-fairness-scored exploitation under dual references), and analysis (no
trajectory classification vs.\ the decay/persistence/collusion-like decision
rule). Its ambivalent reputation effect independently motivates isolating the
contrast we study.

\textbf{Trust, deception, and strategic behavior.} The LH-DECEPTION
benchmark~\cite{lhdeception} models an evolving trust state in a
supervisor--performer hierarchy and audits accumulated deception;
AmongUs~\cite{amongus} elicits emergent strategic misbehavior in a discrete
social-deduction game scored by Elo; StrategicPersuasion~\cite{strategicpersuasion}
formalizes one-shot sender--receiver persuasion with explicit utilities. None
covers symmetric peer bargaining with reputational leverage on both sides
across genuinely repeated interaction: the hierarchy lacks peer symmetry, the
win/lose structure cannot express graded exploitation of a continuous split,
and one-shot games admit no history channel to manipulate.

\textbf{Repeated-game collusion and negotiation foundations.}
Calvano et al.~\cite{calvano2020} show that independent Q-learning agents in
repeated pricing games reliably discover supra-competitive, tacitly collusive
prices without communication; SocialJax~\cite{socialjax} provides sequential
social-dilemma environments with trained MARL baselines exhibiting related
exploitation/cooperation dynamics. Both concern \emph{trained} policies whose
coordination is a property of a learned value table.
Lewis et al.~\cite{lewis2017} established the multi-issue negotiation-dialogue
task our scale-up domain builds on, and Heddaya et al.~\cite{heddaya2023}
collected a human-subjects bargaining corpus contrasting alternating-offer and
free-form negotiation---both as single episodes with no reputation channel.
Training-time approaches such as ARIA~\cite{aria} and
MARSHAL~\cite{marshal} treat strategic games as environments for learning
stronger policies; they measure learning performance, not the behavioral
trajectory of exploitation under a reputation manipulation. None of these
establishes whether dynamics measurable for trained policies also arise for
\emph{prompted}, in-context bargainers conditioned on a natural-language
transcript---the gap our encoding controls (M5) are built to close.

\section{Conclusion}\label{sec:conclusion}

Reputation channels, as deployed for LLM-agent bargaining, condition on
agreement rather than fairness---an asymmetry under which persistent
exploitation and collusion-like coordination are both stable outcomes of the
mechanism assumed to prevent them. We converted the resulting question into a
pre-registered trajectory-classification problem, with dual fairness
references quantified over a tolerance band, a three-conjunct coordination
criterion carrying its own necessity test, and content/order-decoupled
history encodings that firewall the in-context confound no trained-agent
study faces. Across \ph{N=4} backbones and \ph{K=3} domains, reputation
visibility reduced the exploitation trend by \ph{XX\%} yet left
\ph{persistence} as the modal regime in \ph{YY\%} of cells---directionally
corrective, insufficient to self-correct. The governance implication is
stated as a falsifiable conditional: reputation-only safeguards are adequate
only if outcome visibility reliably moves bargaining trajectories toward fair
surplus division, not merely toward stable agreement. Our results indicate it
\ph{does not}, on the backbones and domains tested; the enriched,
fairness-carrying ledger is the natural next intervention, pre-registered
here as a stretch condition. We release the testbed, the full transcript
corpus, and the pre-registration to make every claim auditable.

\balance
\bibliographystyle{IEEEtran}
\bibliography{refs}

\end{document}
